\begin{center}
    \LARGE \textbf{Resumen}
\end{center}


\subsubsection{Contexto}

Estructuración de proyectos de software orientados a la generación automática de aplicaciones que implementen sintetizadores de sonido. Aplicación de ingeniería software al diseño de herramientas sonoras. 

\subsubsection{Descripción del problema}

Cómo estructurar, desde el análisis al diseño, un proyecto software que permita la generación automática de aplicaciones que implementen sintetizadores de sonido, pudiendo ofrecer un mecanismo estandarizado tanto para la especificación de los sintetizadores como para la generación de aplicaciones y su redefinición iterativa.

\subsubsection{Resumen de la solución aportada}

Siguiendo los principios MDA (Model Driven Architecture) se propone una solución basada en la creación de una serie de DSLs (Domain Specific Languages) que guien y acoten las dos principales fases del desarrollo de sintetizadores de sonido en software: la especificación funcional del sintetizador (modelo CIM) y el diseño conceptual del mismo (modelo PIM). Estos modelos estrictos guiados por los DSL desarrollados, permiten acotar el problema, tanto en la generación automática del código fuente como la redefinición iterativa de su análisis y diseño. 

\subsubsection{¿La solución aportada resuelve el problema?}

La solución aportada abre el camino de lo que podría ser resolver el problema, ya que es una propuesta de estandarización más ambiciosa de lo que en este trabajo expongo, en el cual nos enfocaremos más en la generación de código automático que en la redefinición iterativa de los modelos generados.